\section*{Erd\H{o}s Problem 427}

\paragraph{Statement.}
For positive integers \(n,d\), does some \(k\geq1\) satisfy
\[
 d\mid p_{n+1}+\cdots+p_{n+k},
\]
where \(p_j\) is the \(j\)-th prime?

\paragraph{External input.}
We use Shiu's theorem on strings of congruent primes: if \(q\geq3\)
and \(\gcd(a,q)=1\), there exist arbitrarily long strings of consecutive
primes all congruent to \(a\pmod q\).
This is the deep input; its proof is not reproduced here.
It implies that a string of any prescribed finite length exists beyond
any prescribed index: otherwise strings of unbounded length would have
bounded starting index and would force all sufficiently late primes to
lie in that residue class, which also supplies arbitrarily late strings.

\paragraph{Deduction of the requested divisibility.}
The case \(d=1\) is immediate. For \(d=2\), the primes after \(p_n\)
are odd because \(n\geq1\); two of them have an even sum.
For \(d\geq3\), apply Shiu's theorem with \(a=1\), \(q=d\).
Choose \(m>n\) such that
\[
 p_m,\ldots,p_{m+d-1}\equiv1\pmod d.
\]
Write \(s=\sum_{j=n+1}^{m-1}p_j\), allowing an empty sum. There is a
unique \(\ell\in\{1,\ldots,d\}\) with \(\ell\equiv-s\pmod d\).
Then
\[
 \sum_{j=n+1}^{m+\ell-1}p_j
 \equiv s+\ell\equiv0\pmod d.
\]
Taking \(k=m+\ell-1-n\geq1\) proves the assertion. Choosing strings
arbitrarily far out gives infinitely many such \(k\) for each fixed
\(n,d\).

\paragraph{Attribution and scope.}
This reconstructs the deduction attributed to Cedric Pilatte on the
problem page. Its external input is
D.~K.~L.~Shiu, \emph{Strings of congruent primes},
Journal of the London Mathematical Society 61 (2000), 359--373,
\url{https://doi.org/10.1112/S0024610799007863}.
The current statement is \url{https://www.erdosproblems.com/427},
checked 24 September 2026. Completeness is relative to the explicitly
stated established theorem, not an independent proof of Shiu's theorem.

\paragraph{Completion estimate.}
The full assertion follows from Shiu's theorem with all quantifiers checked.
\textbf{COMPLETION: 100\%}
